% Macros for reviewer responses and back-links from response document
% to manuscript line numbers. Pattern adapted from popgensbi manuscript.
%
% Expects the lineno package to be loaded by the calling document and
% \linenumbers to be active in the manuscript body.

% Reviewer / point counters
\newcommand{\thereviewer}{}
\newcounter{point}
\setcounter{point}{0}

% Start a new reviewer block. Argument is the reviewer label
% (e.g. \reviewersection{1} or \reviewersection{AE}).
\newcommand{\reviewersection}[1]{%
  \renewcommand{\thereviewer}{#1}%
  \setcounter{point}{0}%
  \section*{Reviewer \thereviewer:}%
}

% A single review point. Argument is an optional name (often empty).
% Pattern from
% http://tex.stackexchange.com/questions/2317/latex-style-or-macro-for-detailed-response-to-referee-report
\newenvironment{point}[1]
  {\color{gray}\refstepcounter{point}\bigskip\hrule\medskip\noindent
   \slshape{\fontseries{b}\selectfont (\thereviewer.\thepoint) #1}}
  {}

\newcommand{\reply}{\normalfont\medskip\noindent\textbf{Reply}:\ }

% \revpoint{R}{P} marks a manuscript line that addresses reviewer R, point P.
% Use this in the manuscript wherever you make a change in response to a
% reviewer point so the response document can back-link to it.
\newcommand{\revpoint}[2]{\hypertarget{llineno:rev#1:#2}{\linelabel{rr:rev#1:#2}}}

% \revreffull{R}{P} renders a (page, line) reference to such a target.
\newcommand{\revreffull}[2]{{(p.\ \hyperlink{llineno:rev#1:#2}{\pageref{rr:rev#1:#2}, l.\ \lineref{rr:rev#1:#2}})}}

% \revref expands to the page/line reference for the current reviewer/point
% inside a {point} environment.
\newcommand{\revref}{\revreffull{\thereviewer}{\thepoint}\xspace}

% Named back-links: place \llabel{name} in the manuscript, refer to it
% from anywhere with \llname{name}.
\newcommand{\llabel}[1]{\hypertarget{ll:#1}{\linelabel{#1}}}
\newcommand{\llname}[1]{{(p.\ \hyperlink{ll:#1}{\pageref{#1}, l.\ \lineref{#1}})}}

% Italics-quoted block (used to set off the verbatim reviewer text).
\newenvironment{itquote}
  {\begin{quote}\color{gray}\itshape}
  {\end{quote}}

% Convenience macros
\newcommand{\rollover}{\reply{The reviewer makes an excellent point that we have missed out entirely. We have made all the changes suggested, down to the minutiae \revref.}}
\newcommand{\playdead}{\reply{The reviewer makes an excellent point. We have made an utterly trivial change \revref{}that we think deals entirely with the concern raised.}}

% lineno + amsmath fix: lineno doesn't number lines inside align/equation
% environments by default. Patch from
% http://tex.stackexchange.com/questions/43648
\ifcsname patchAmsMathEnvironmentForLineno\endcsname
\else
\newcommand*\patchAmsMathEnvironmentForLineno[1]{%
  \expandafter\let\csname old#1\expandafter\endcsname\csname #1\endcsname
  \expandafter\let\csname oldend#1\expandafter\endcsname\csname end#1\endcsname
  \renewenvironment{#1}%
    {\linenomath\csname old#1\endcsname}%
    {\csname oldend#1\endcsname\endlinenomath}}
\newcommand*\patchBothAmsMathEnvironmentsForLineno[1]{%
  \patchAmsMathEnvironmentForLineno{#1}%
  \patchAmsMathEnvironmentForLineno{#1*}}
\AtBeginDocument{%
  \patchBothAmsMathEnvironmentsForLineno{equation}%
  \patchBothAmsMathEnvironmentsForLineno{align}%
  \patchBothAmsMathEnvironmentsForLineno{flalign}%
  \patchBothAmsMathEnvironmentsForLineno{alignat}%
  \patchBothAmsMathEnvironmentsForLineno{gather}%
  \patchBothAmsMathEnvironmentsForLineno{multline}%
}
\fi
